%% toc: P3 30 30.1,30.2,30.3,30.4
\section{Substructure, extension and chain preservation: proofs}
\label{sec:annex-p03-w03-p03-026}

\subsection{Modern preservation arguments}
\label{sec:annex-p03-w03-p03-026-modern}

The items below prove the results of
\cref{sec:p03-w03-p03-026-preservation-proofs}.  Compactness is used in the
form of \cref{thm:fol-compactness}; the diagram tests, isomorphism
invariance and the elementary chain theorem are
\cref{thm:diagram-tests,lem:iso-invariance,thm:elementary-chain}.
Two further facts are used repeatedly.  \emph{Expansion by new constants}:
if constants \(\bar c\) occur in no member of a set \(\Sigma\) of sentences,
\(\theta(\bar y)\) is any formula without those constants, and
\(\Sigma\entails\neg\theta(\bar c)\), then
\(\Sigma\entails\forall\bar y\,\neg\theta(\bar y)\); for, given a model
\(\mathcal N\) of \(\Sigma\) and a tuple \(\bar b\) from \(N\), the expansion
of \(\mathcal N\) interpreting \(\bar c\) by \(\bar b\) still satisfies
\(\Sigma\) by \cref{lem:reduct-invariance}, hence satisfies
\(\neg\theta(\bar c)\), so \(\mathcal N\entails\neg\theta(\bar b)\).
\emph{Copying}: if \(h:\mathcal A\to\mathcal D\) is an embedding of
\(\Lang\)-structures, there are \(\mathcal B\supseteq\mathcal A\) and an
isomorphism \(g:\mathcal B\to\mathcal D\) extending \(h\).  Put
\(X=(D\setminus h[A])\times\{A\}\), with Kuratowski pairs
\((d,A)=\{\{d\},\{d,A\}\}\).  No pair \((d,A)\) belongs to \(A\), for
otherwise \(A\in\{d,A\}\in(d,A)\in A\) would be a membership cycle, which
the axiom of foundation excludes; so \(X\)
is disjoint from \(A\).  Let \(g\) be \(h\) on \(A\) and \((d,A)\mapsto d\)
on \(X\).  Then \(g\) is a bijection from \(A\cup X\) onto \(D\), because
\(h\) is injective.  Transport the structure: \(R^{\mathcal B}(\bar e)\) iff
\(R^{\mathcal D}(g\bar e)\), \(f^{\mathcal B}(\bar e)=g^{-1}(f^{\mathcal
D}(g\bar e))\) and \(c^{\mathcal B}=g^{-1}(c^{\mathcal D})\).  Then \(g\) is
an isomorphism, and for \(\bar a\) from \(A\) we get
\(f^{\mathcal B}(\bar a)=g^{-1}(h(f^{\mathcal A}(\bar a)))=f^{\mathcal A}(\bar a)\)
and \(R^{\mathcal B}(\bar a)\iff R^{\mathcal D}(h\bar a)\iff R^{\mathcal A}(\bar a)\);
likewise for constants.  Thus \(\mathcal A\subseteq\mathcal B\).

\begin{annexitem}{annex:p03-w03-p03-026-union}{thm:p03-w03-p03-026-union}{The union of a chain}
\begin{proof}
Let \(M=\bigcup_{i<\delta}\dom\mathcal M_i\); it is nonempty because
\(\delta>0\) and \(\dom\mathcal M_0\neq\varnothing\).  A constant symbol
\(c\) has \(c^{\mathcal M_i}=c^{\mathcal M_0}\) for all \(i\), because
\(\mathcal M_0\subseteq\mathcal M_i\); put \(c^{\mathcal M}=c^{\mathcal M_0}\).
For an \(n\)-ary operation symbol \(f\) and \(\bar a\in M^n\), each
\(a_k\) lies in some \(\dom\mathcal M_{i_k}\); with \(i=\max_k i_k\), the
whole tuple lies in \(\dom\mathcal M_i\).  Put
\(f^{\mathcal M}(\bar a)=f^{\mathcal M_i}(\bar a)\).  If also
\(\bar a\in(\dom\mathcal M_j)^n\) and, say, \(i\le j\), then
\(f^{\mathcal M_j}(\bar a)=f^{\mathcal M_i}(\bar a)\) because
\(\mathcal M_i\subseteq\mathcal M_j\); the value does not depend on the
stage.  For an \(n\)-ary relation symbol put
\(R^{\mathcal M}=\bigcup_{i<\delta}R^{\mathcal M_i}\).

Fix \(i\).  The domain of \(\mathcal M_i\) is contained in \(M\);
constants and operations of \(\mathcal M\) agree with those of
\(\mathcal M_i\) on \(\dom\mathcal M_i\) by construction.  For relations let
\(\bar a\in(\dom\mathcal M_i)^n\).  If \(\bar a\in R^{\mathcal M_i}\), then
\(\bar a\in R^{\mathcal M}\).  Conversely, if \(\bar a\in R^{\mathcal M}\),
then \(\bar a\in R^{\mathcal M_j}\) for some \(j\).  If \(j\le i\), then
\(R^{\mathcal M_j}\subseteq R^{\mathcal M_i}\) since
\(\mathcal M_j\subseteq\mathcal M_i\).  If \(i<j\), then
\(R^{\mathcal M_i}=R^{\mathcal M_j}\cap(\dom\mathcal M_i)^n\) since
\(\mathcal M_i\subseteq\mathcal M_j\), and \(\bar a\) lies in this
intersection.  Either way \(\bar a\in R^{\mathcal M_i}\).  Hence
\(R^{\mathcal M}\cap(\dom\mathcal M_i)^n=R^{\mathcal M_i}\), and
\(\mathcal M_i\subseteq\mathcal M\).
\end{proof}
\end{annexitem}

\begin{annexitem}{annex:p03-w03-p03-026-easy}{thm:p03-w03-p03-026-easy}{Preservation by the syntactic sides}
\begin{proof}
(1) By induction on terms, every term \(t(\bar x)\) has
\(t^{\mathcal M}(\bar a)=t^{\mathcal N}(\bar a)\) for \(\bar a\) from \(M\):
variables and constants are immediate, and an application
\(f(t_1,\ldots,t_n)\) is evaluated by \(f^{\mathcal M}\), which is the
restriction of \(f^{\mathcal N}\).  An equality atom \(t_1=t_2\) is
therefore true in \(\mathcal M\) iff it is true in \(\mathcal N\), equality
being identity in both.  A relational atom \(R(t_1,\ldots,t_n)\) has the
same arguments in both structures, and
\(R^{\mathcal M}=R^{\mathcal N}\cap M^n\).  The clauses for
\(\neg,\land,\lor,\to\) depend only on the truth values of the parts, so the
equivalence passes to every quantifier-free formula.

(2) Suppose \(\mathcal N\entails\forall\bar x\,\theta(\bar x,\bar a)\).
Every tuple \(\bar b\) from \(M\) is a tuple from \(N\), so
\(\mathcal N\entails\theta(\bar b,\bar a)\) and, by (1),
\(\mathcal M\entails\theta(\bar b,\bar a)\).  Hence
\(\mathcal M\entails\forall\bar x\,\theta(\bar x,\bar a)\).  If
\(\mathcal M\entails\exists\bar x\,\theta(\bar x,\bar a)\) with witness
\(\bar b\) from \(M\), then by (1) the same \(\bar b\) witnesses
\(\mathcal N\entails\exists\bar x\,\theta(\bar x,\bar a)\).

(3) Let \(\mathcal M\) be the union of a chain
\(\langle\mathcal M_i:i<\delta\rangle\) whose members satisfy
\(\forall\bar x\,\exists\bar y\,\theta(\bar x,\bar y)\).  Let \(\bar a\) be
a tuple from \(M\).  Its finitely many entries lie in stages
\(i_1,\ldots,i_n\); with \(i\) their maximum (or \(i=0\) when \(\bar x\) is
empty), \(\bar a\) lies in \(\dom\mathcal M_i\).  Then
\(\mathcal M_i\entails\exists\bar y\,\theta(\bar a,\bar y)\); choose a
witness \(\bar b\) from \(\dom\mathcal M_i\).  Since
\(\mathcal M_i\subseteq\mathcal M\) by \cref{thm:p03-w03-p03-026-union},
part (1) gives \(\mathcal M\entails\theta(\bar a,\bar b)\).  As \(\bar a\) was
arbitrary, \(\mathcal M\entails\forall\bar x\,\exists\bar y\,\theta\).
\end{proof}
\end{annexitem}

\begin{annexitem}{annex:p03-w03-p03-026-prenex}{thm:p03-w03-p03-026-prenex}{Negations, conjunctions and disjunctions in the prefix classes}
\begin{proof}
By the satisfaction clauses, a structure satisfies
\(\neg\forall\bar x\,\theta\) iff some tuple falsifies \(\theta\), iff it
satisfies \(\exists\bar x\,\neg\theta\); and \(\neg\theta\) is
quantifier-free.  Symmetrically \(\neg\exists\bar x\,\theta\) is equivalent
to \(\forall\bar x\,\neg\theta\).

For two universal sentences \(\forall\bar x\,\theta_1\) and
\(\forall\bar y\,\theta_2\), first rename bound variables so that
\(\bar x\) and \(\bar y\) are disjoint; renaming bound variables does not
change truth (\cref{cor:alpha-invariance}).  We claim
\[
 \forall\bar x\,\theta_1\land\forall\bar y\,\theta_2
 \quad\text{is equivalent to}\quad
 \forall\bar x\,\forall\bar y\,(\theta_1\land\theta_2).
\]
If the left side holds, every pair of tuples \(\bar a,\bar b\) satisfies
both conjuncts, since \(\theta_1\) has no free \(\bar y\) and \(\theta_2\)
no free \(\bar x\).  If the right side holds and \(\bar a\) is given, choose
any tuple \(\bar b\), which exists because domains are nonempty; then
\(\theta_1(\bar a)\) holds.  The second conjunct is obtained in the same
way.  Induction on the number of conjuncts handles finite conjunctions; the
empty conjunction is \(\forall x\,(x=x)\) by convention.  For existential
sentences,
\(\exists\bar x\,\theta_1\lor\exists\bar y\,\theta_2\) is equivalent to
\(\exists\bar x\,\exists\bar y\,(\theta_1\lor\theta_2)\): a witness for
either disjunct together with an arbitrary tuple for the other block
witnesses the right side, and a witness \((\bar a,\bar b)\) for the right
side witnesses one of the disjuncts.  Parameters are constants and do not
interact with these steps.
\end{proof}
\end{annexitem}

\begin{annexitem}{annex:p03-w03-p03-026-embed}{thm:p03-w03-p03-026-embed}{Universal consequences and embeddings into models}
\begin{proof}
(2)\(\Rightarrow\)(1).  Let \(\mathcal M\subseteq\mathcal N\entails T\) and
\(\psi\in T_{\forall}\).  Then \(\mathcal N\entails\psi\), and
\(\mathcal M\entails\psi\) by \cref{thm:p03-w03-p03-026-easy}(2).

(1)\(\Rightarrow\)(2).  Assume \(\mathcal M\entails T_{\forall}\) and put
\(\Sigma=T\cup\operatorname{Diag}(\mathcal M)\), a set of
\(\Lang(M)\)-sentences.  Suppose \(\Sigma\) has no model.  By compactness
some finite subset has no model; enlarging it, we may take it to be
\(T_0\cup\{\delta_1,\ldots,\delta_k\}\) with \(T_0\subseteq T\) finite and
\(\delta_j\in\operatorname{Diag}(\mathcal M)\). If \(k=0\), then \(T_0\)
has no model, so \(\forall x\,\neg(x=x)\in T_\forall\), contradicting
\(\mathcal M\entails T_\forall\). Thus assume \(k>0\).
Let \(\bar a\) list without
repetition the elements of \(M\) whose constants occur in the
\(\delta_j\), and write
\(\delta_1\land\cdots\land\delta_k=\delta(\bar c_{\bar a})\) with
\(\delta(\bar x)\) a quantifier-free \(\Lang\)-formula.  Then \(T\entails\neg\delta(\bar c_{\bar a})\).
The constants \(\bar c_{\bar a}\) do not occur in \(T\), so expansion by new
constants gives \(T\entails\forall\bar x\,\neg\delta(\bar x)\).  This is a
universal sentence, so it belongs to \(T_{\forall}\) and holds in
\(\mathcal M\).  But every \(\delta_j\) is true in \(\mathcal M_M\), so
\(\mathcal M\entails\delta(\bar a)\), a contradiction.

Let \(\mathcal N'\) be a model of \(\Sigma\).  The map
\(h:a\mapsto c_a^{\mathcal N'}\) is an embedding of \(\mathcal M\) into the
\(\Lang\)-reduct of \(\mathcal N'\), as in the proof of
\cref{thm:diagram-tests}(1): \(\operatorname{Diag}(\mathcal M)\) contains
\(\neg(c_a=c_b)\) for \(a\neq b\), \(f(c_{a_1},\ldots,c_{a_n})=c_{f(\bar a)}\),
\(c=c_{c^{\mathcal M}}\), and each true atomic or negated atomic relation
sentence, so \(h\) is injective and preserves operations, constants and
relations in both directions.  Copying gives \(\mathcal N\supseteq\mathcal M\) and an
isomorphism \(g:\mathcal N\to\mathcal N'{\restriction}\Lang\) extending
\(h\).  The reduct satisfies \(T\) (\cref{lem:reduct-invariance}), so
\(\mathcal N\entails T\) by \cref{lem:iso-invariance}.
\end{proof}
\end{annexitem}

\begin{annexitem}{annex:p03-w03-p03-026-los-tarski-relative}{thm:p03-w03-p03-026-los-tarski-relative}{The relative \L o\'s--Tarski theorem}
\begin{proof}
(2)\(\Rightarrow\)(1).  Let \(\mathcal M\entails S\) and
\(\mathcal M\subseteq\mathcal N\entails T\).  Since
\(\operatorname{Mod}(T)=\operatorname{Mod}(S\cup U)\), \(\mathcal N\)
satisfies every member of \(U\), and each of these universal sentences
passes down to \(\mathcal M\) by \cref{thm:p03-w03-p03-026-easy}(2).  Thus
\(\mathcal M\entails S\cup U\), so \(\mathcal M\entails T\).

(1)\(\Rightarrow\)(2), with \(U=T_{\forall}\).  Since \(S\subseteq T\) and
\(T\entails T_{\forall}\), every model of \(T\) is a model of
\(S\cup T_{\forall}\).  Conversely let \(\mathcal M\entails S\cup T_{\forall}\).
By \cref{thm:p03-w03-p03-026-embed}, \(\mathcal M\subseteq\mathcal N\) for
some \(\mathcal N\entails T\).  As \(\mathcal M\entails S\), hypothesis (1)
gives \(\mathcal M\entails T\).  Hence
\(\operatorname{Mod}(T)=\operatorname{Mod}(S\cup T_{\forall})\).
\end{proof}
\end{annexitem}

\begin{annexitem}{annex:p03-w03-p03-026-los-tarski-class}{thm:p03-w03-p03-026-los-tarski-class}{The absolute class form}
\begin{proof}
Let \(\mathcal K=\operatorname{Mod}(T)\).  If \(\mathcal K\) is closed under
substructures, then hypothesis (1) of
\cref{thm:p03-w03-p03-026-los-tarski-relative} holds with
\(S=\varnothing\), because every structure is a model of \(\varnothing\) and
every substructure of a member of \(\mathcal K\) is in \(\mathcal K\).  The
theorem gives \(\mathcal K=\operatorname{Mod}(T_{\forall})\), an
axiomatization by universal sentences.  Conversely, if
\(\mathcal K=\operatorname{Mod}(U)\) with \(U\) universal,
\(\mathcal N\in\mathcal K\) and \(\mathcal M\subseteq\mathcal N\), then each
member of \(U\) passes from \(\mathcal N\) to \(\mathcal M\) by
\cref{thm:p03-w03-p03-026-easy}(2), so \(\mathcal M\in\mathcal K\).
\end{proof}
\end{annexitem}

\begin{annexitem}{annex:p03-w03-p03-026-los-tarski-sentence}{thm:p03-w03-p03-026-los-tarski-sentence}{The sentence form and its dual}
\begin{proof}
Item 1, (b)\(\Rightarrow\)(a).  Let \(\mathcal M\subseteq\mathcal N\) be
models of \(S\) with \(\mathcal M\entails\varphi\).  Then
\(\mathcal M\entails\psi\), the existential sentence \(\psi\) passes to
\(\mathcal N\) by \cref{thm:p03-w03-p03-026-easy}(2), and
\(\mathcal N\entails\varphi\) because \(\mathcal N\entails S\).

Item 1, (a)\(\Rightarrow\)(b).  Put \(T=S\cup\{\neg\varphi\}\).  We check
hypothesis (1) of \cref{thm:p03-w03-p03-026-los-tarski-relative}.  Let
\(\mathcal M\entails S\) with \(\mathcal M\subseteq\mathcal N\entails T\).
If \(\mathcal M\entails\varphi\), then (a), applied to the models
\(\mathcal M\subseteq\mathcal N\) of \(S\), gives
\(\mathcal N\entails\varphi\), contradicting \(\mathcal N\entails\neg\varphi\).
So \(\mathcal M\entails\neg\varphi\) and \(\mathcal M\entails T\).  The
theorem yields \(\operatorname{Mod}(T)=\operatorname{Mod}(S\cup T_{\forall})\),
so \(S\cup T_{\forall}\entails\neg\varphi\).  By compactness there are finite
\(S_0\subseteq S\) and \(U_0\subseteq T_{\forall}\) with
\(S_0\cup U_0\entails\neg\varphi\).  By
\cref{thm:p03-w03-p03-026-prenex} the conjunction of \(U_0\) is equivalent to
a universal sentence \(\chi\) (the sentence \(\forall x\,(x=x)\) if
\(U_0=\varnothing\)).  Then \(S\entails\chi\to\neg\varphi\).  Each member of
\(U_0\) is a consequence of \(T\), so \(S\cup\{\neg\varphi\}\entails\chi\),
that is, \(S\entails\neg\varphi\to\chi\).  Hence
\(S\entails\varphi\leftrightarrow\neg\chi\), and \(\neg\chi\) is equivalent
to an existential sentence \(\psi\) by \cref{thm:p03-w03-p03-026-prenex}.

Item 2.  Condition 2(a) for \(\varphi\) is the contrapositive of condition
1(a) for \(\neg\varphi\): both say that for models
\(\mathcal M\subseteq\mathcal N\) of \(S\), \(\mathcal M\entails\neg\varphi\)
implies \(\mathcal N\entails\neg\varphi\).  By item 1 this holds iff
\(S\entails\neg\varphi\leftrightarrow\psi\) for some existential \(\psi\).
If so, \(\neg\psi\) is equivalent to a universal \(\chi\) and
\(S\entails\varphi\leftrightarrow\chi\).  Conversely, if
\(S\entails\varphi\leftrightarrow\chi\) with \(\chi\) universal, then
\(\neg\chi\) is equivalent to an existential \(\psi\) and
\(S\entails\neg\varphi\leftrightarrow\psi\).  Thus 2(a) and 2(b) are
equivalent.
\end{proof}
\end{annexitem}

\begin{annexitem}{annex:p03-w03-p03-026-step-up}{thm:p03-w03-p03-026-step-up}{Finding a model of \(T\) above \(\mathcal A\)}
\begin{proof}
Let \(\Sigma=T\cup\operatorname{Diag}_{\forall}(\mathcal A)\), a set of
\(\Lang(A)\)-sentences.  Suppose \(\Sigma\) has no model.  By compactness a
finite subset \(T_0\cup\{\chi_1,\ldots,\chi_k\}\) has no model, with
\(T_0\subseteq T\) and \(\chi_j\in\operatorname{Diag}_{\forall}(\mathcal A)\).
By \cref{thm:p03-w03-p03-026-prenex}, \(\chi_1\land\cdots\land\chi_k\) is
equivalent to one universal sentence
\(\forall\bar x\,\theta(\bar x,\bar c_{\bar a})\), with
\(\theta(\bar x,\bar z)\) a quantifier-free \(\Lang\)-formula and \(\bar a\)
a tuple of distinct elements of \(A\); this sentence is true in
\(\mathcal A_A\) because each \(\chi_j\) is. When \(k=0\), take
\(\bar x=(x)\), \(\bar z\) empty and \(\theta(x)=x=x\), so this is
the universal sentence \(\forall x\,(x=x)\).  Then
\(T\entails\neg\forall\bar x\,\theta(\bar x,\bar c_{\bar a})\), that is,
\(T\entails\exists\bar x\,\neg\theta(\bar x,\bar c_{\bar a})\).  The
constants \(\bar c_{\bar a}\) do not occur in \(T\); expansion by new
constants, applied to the formula \(\forall\bar x\,\theta(\bar x,\bar z)\),
gives \(T\entails\forall\bar z\,\neg\forall\bar x\,\theta(\bar x,\bar z)\),
that is,
\[
 T\entails\forall\bar z\,\exists\bar x\,\neg\theta(\bar x,\bar z).
\]
This is a \(\forall\exists\) sentence, so it belongs to
\(T_{\forall\exists}\) and holds in \(\mathcal A\).  In particular
\(\mathcal A\entails\exists\bar x\,\neg\theta(\bar x,\bar a)\), contradicting
\(\mathcal A_A\entails\forall\bar x\,\theta(\bar x,\bar c_{\bar a})\).

So \(\Sigma\) has a model \(\mathcal B'\).  It satisfies
\(\operatorname{Diag}(\mathcal A)\subseteq\operatorname{Diag}_{\forall}(\mathcal A)\),
so, by the same argument as in \cref{annex:p03-w03-p03-026-embed} (the one
in the proof of \cref{thm:diagram-tests}(1)), the map
\(h:a\mapsto c_a^{\mathcal B'}\) embeds \(\mathcal A\) into the
\(\Lang\)-reduct of \(\mathcal B'\).  Copying gives
\(\mathcal B\supseteq\mathcal A\) and an isomorphism
\(g:\mathcal B\to\mathcal B'{\restriction}\Lang\) extending \(h\).  The
expansion \(\mathcal B_A\) interpreting each \(c_a\) by \(a\) is isomorphic
to \(\mathcal B'\) via \(g\), since \(g(a)=c_a^{\mathcal B'}\).  By
\cref{lem:iso-invariance,lem:reduct-invariance},
\(\mathcal B\entails T\) and
\(\mathcal B_A\entails\operatorname{Diag}_{\forall}(\mathcal A)\).

Finally let \(\varepsilon\) be an existential \(\Lang(A)\)-sentence true in
\(\mathcal B_A\).  If \(\varepsilon\) were false in \(\mathcal A_A\), then
\(\neg\varepsilon\), which is equivalent to a universal
\(\Lang(A)\)-sentence \(\varepsilon'\) by \cref{thm:p03-w03-p03-026-prenex},
would be true in \(\mathcal A_A\).  Then
\(\varepsilon'\in\operatorname{Diag}_{\forall}(\mathcal A)\) holds in
\(\mathcal B_A\), so \(\mathcal B_A\entails\neg\varepsilon\), a
contradiction.
\end{proof}
\end{annexitem}

\begin{annexitem}{annex:p03-w03-p03-026-step-back}{thm:p03-w03-p03-026-step-back}{Amalgamating \(\mathcal B\) with the elementary diagram of \(\mathcal A\)}
\begin{proof}
Work in \(\Lang(B)\), which contains \(c_a\) for every \(a\in A\subseteq B\).
Let \(\Sigma=\operatorname{ElDiag}(\mathcal A)\cup\operatorname{Diag}(\mathcal B)\).
Suppose \(\Sigma\) has no model.  By compactness some finite
\(E_0\cup\{\delta_1,\ldots,\delta_k\}\) has no model, with
\(E_0\subseteq\operatorname{ElDiag}(\mathcal A)\) and
\(\delta_j\in\operatorname{Diag}(\mathcal B)\). The case \(k=0\) is
impossible, since \(\mathcal A_A\) satisfies \(E_0\). Write
\(\delta_1\land\cdots\land\delta_k=\delta(\bar c_{\bar a},\bar c_{\bar b})\),
where \(\bar a\) lists distinct elements of \(A\), \(\bar b\) lists distinct
elements of \(B\setminus A\), and \(\delta(\bar x,\bar y)\) is a
quantifier-free \(\Lang\)-formula.
Then \(\operatorname{ElDiag}(\mathcal A)\entails
\neg\delta(\bar c_{\bar a},\bar c_{\bar b})\).  The constants
\(\bar c_{\bar b}\) do not occur in \(\operatorname{ElDiag}(\mathcal A)\), so
by expansion by new constants
\(\operatorname{ElDiag}(\mathcal A)\entails
\forall\bar y\,\neg\delta(\bar c_{\bar a},\bar y)\).  Since
\(\mathcal A_A\entails\operatorname{ElDiag}(\mathcal A)\),
\[
 \mathcal A_A\entails\neg\exists\bar y\,\delta(\bar c_{\bar a},\bar y).
\]
But each \(\delta_j\) is true in \(\mathcal B_B\), so \(\bar b\) witnesses
\(\mathcal B_A\entails\exists\bar y\,\delta(\bar c_{\bar a},\bar y)\).  This
existential \(\Lang(A)\)-sentence is then true in \(\mathcal A_A\) by
hypothesis, a contradiction.

Let \(\mathcal C'\) be a model of \(\Sigma\).  As in \cref{annex:p03-w03-p03-026-embed}, the map
\(g:b\mapsto c_b^{\mathcal C'}\) embeds \(\mathcal B\) into
\(\mathcal C'{\restriction}\Lang\), because
\(\mathcal C'\entails\operatorname{Diag}(\mathcal B)\).  Let
\(f=g{\restriction}A\), so \(f(a)=c_a^{\mathcal C'}\).  If
\(\mathcal A\entails\varphi(\bar a)\), then \(\varphi(\bar c_{\bar a})\) lies
in \(\operatorname{ElDiag}(\mathcal A)\), so
\(\mathcal C'{\restriction}\Lang\entails\varphi(f\bar a)\) by
\cref{lem:reduct-invariance}; if \(\mathcal A\entails\neg\varphi(\bar a)\),
the same argument applies to \(\neg\varphi\).  Hence \(f\) is an elementary
embedding of \(\mathcal A\) into \(\mathcal C'{\restriction}\Lang\), the
direction used in the proof of \cref{thm:diagram-tests}(2).  Copying gives \(\mathcal C\supseteq\mathcal B\) and
an isomorphism \(k:\mathcal C\to\mathcal C'{\restriction}\Lang\) extending
\(g\).  Then \(\mathcal A\subseteq\mathcal B\subseteq\mathcal C\).  For a
formula \(\varphi(\bar x)\) and \(\bar a\) from \(A\),
\[
 \mathcal A\entails\varphi(\bar a)
 \iff\mathcal C'{\restriction}\Lang\entails\varphi(f\bar a)
 \iff\mathcal C\entails\varphi(k^{-1}f\bar a)
 \iff\mathcal C\entails\varphi(\bar a),
\]
by elementarity of \(f\), \cref{lem:iso-invariance}, and
\(k^{-1}(f(a))=k^{-1}(g(a))=a\).  Hence \(\mathcal A\preccurlyeq\mathcal C\).
\end{proof}
\end{annexitem}

\begin{annexitem}{annex:p03-w03-p03-026-chang-los-suszko}{thm:p03-w03-p03-026-chang-los-suszko}{The Chang--\L o\'s--Suszko theorem: all four directions}
\begin{proof}
(4)\(\Rightarrow\)(1).  Let \(\operatorname{Mod}(T)=\operatorname{Mod}(\Gamma)\)
with \(\Gamma\) a set of \(\forall\exists\) sentences.  Every member of a
chain of models of \(T\) satisfies \(\Gamma\), so the union does by
\cref{thm:p03-w03-p03-026-easy}(3); hence it is a model of \(T\).

(1)\(\Rightarrow\)(2).  A chain indexed by \(\omega\) is a chain in the sense
of \cref{def:p03-w03-p03-026-chain}.

(2)\(\Rightarrow\)(3).  Since \(T\entails T_{\forall\exists}\), every model
of \(T\) is a model of \(T_{\forall\exists}\).  Conversely let
\(\mathcal M\entails T_{\forall\exists}\).  Define \(\mathcal A_k\) and
\(\mathcal B_k\) for \(k<\omega\) by recursion, keeping
\(\mathcal A_k\entails T_{\forall\exists}\).  Put
\(\mathcal A_0=\mathcal M\).  Given \(\mathcal A_k\), choose by
\cref{thm:p03-w03-p03-026-step-up} a model \(\mathcal B_k\) of \(T\) with
\(\mathcal A_k\subseteq\mathcal B_k\) such that existential
\(\Lang(A_k)\)-sentences true in \(\mathcal B_k\) are true in
\(\mathcal A_k\).  Then choose by \cref{thm:p03-w03-p03-026-step-back} a
structure \(\mathcal A_{k+1}\) with
\(\mathcal A_k\preccurlyeq\mathcal A_{k+1}\) and
\(\mathcal B_k\subseteq\mathcal A_{k+1}\).  Elementary extensions satisfy
the same sentences, so \(\mathcal A_{k+1}\entails T_{\forall\exists}\).
(The recursion selects a structure at each stage, and compactness is applied
to signatures \(\Lang(A_k)\) that are uncountable when \(A_k\) is; the
construction therefore uses the Axiom of Choice.)

The sequence
\(\mathcal A_0\subseteq\mathcal B_0\subseteq\mathcal A_1\subseteq\mathcal B_1\subseteq\cdots\)
is a chain indexed by \(\omega\), because inclusion of substructures is
transitive.  Let \(\mathcal C\) be its union
(\cref{thm:p03-w03-p03-026-union}).  The subsequences have the same union:
\(\bigcup_k\dom\mathcal A_k=\bigcup_k\dom\mathcal B_k=\dom\mathcal C\)
because \(A_k\subseteq B_k\subseteq A_{k+1}\); each relation of
\(\mathcal C\) is
\(\bigcup_kR^{\mathcal A_k}\cup\bigcup_kR^{\mathcal B_k}\), and
\(R^{\mathcal A_k}\subseteq R^{\mathcal B_k}\subseteq R^{\mathcal A_{k+1}}\),
so it equals both \(\bigcup_kR^{\mathcal A_k}\) and \(\bigcup_kR^{\mathcal B_k}\);
operations and constants are computed at any stage containing the arguments.
Hence \(\mathcal C=\bigcup_k\mathcal B_k=\bigcup_k\mathcal A_k\).

The \(\mathcal B_k\) form a chain of models of \(T\) indexed by \(\omega\),
since \(\mathcal B_k\subseteq\mathcal A_{k+1}\subseteq\mathcal B_{k+1}\).
By (2), \(\mathcal C\entails T\).  The \(\mathcal A_k\) form an elementary
chain: for \(k<l\), \(\mathcal A_k\preccurlyeq\mathcal A_l\) because
elementary inclusions compose.  As \(\omega\) is a limit ordinal, the
elementary chain theorem gives \(\mathcal A_0\preccurlyeq\mathcal C\).
Therefore \(\mathcal M=\mathcal A_0\entails T\).  This proves
\(\operatorname{Mod}(T_{\forall\exists})\subseteq\operatorname{Mod}(T)\).

(3)\(\Rightarrow\)(4).  \(T_{\forall\exists}\) consists of
\(\forall\exists\) sentences.
\end{proof}
\end{annexitem}

\begin{annexitem}{annex:p03-w03-p03-026-elementary-chains-differ}{thm:p03-w03-p03-026-elementary-chains-differ}{Elementary chains against arbitrary chains}
\begin{proof}
Let \(\mathcal K=\operatorname{Mod}(T)\) and let
\(\langle\mathcal M_i:i<\delta\rangle\) be an elementary chain of members
of \(\mathcal K\) with union \(\mathcal M\).  If \(\delta\) has a largest
element \(i\), then \(\mathcal M=\mathcal M_i\in\mathcal K\).  Otherwise
\(\delta\) is a limit ordinal, \(\mathcal M_0\preccurlyeq\mathcal M\) by
\cref{thm:elementary-chain}, and \(\mathcal M\entails T\) because
\(\mathcal M_0\) does.

In the signature with one binary relation symbol \(\le\), the class of
linear orders with a greatest element is \(\operatorname{Mod}(T)\) for
\(T\) consisting of the reflexivity, antisymmetry, transitivity and
totality axioms together with
\(\exists x\,\forall y\,(y\le x)\).  For \(n<\omega\) let
\(\mathcal M_n=(\{0,\ldots,n\},\le)\) with the usual order.  Each is a
substructure of the next, since its order is the restriction, and each has
greatest element \(n\).  The union is \((\N,\le)\), which has no greatest
element.  So \(\operatorname{Mod}(T)\) is not closed under unions of chains
of substructures, and by the implication (4)\(\Rightarrow\)(1) of
\cref{thm:p03-w03-p03-026-chang-los-suszko} it has no \(\forall\exists\)
axiomatization.  By the first paragraph it is nevertheless closed under
unions of elementary chains.
\end{proof}
\end{annexitem}

\begin{annexitem}{annex:p03-w03-p03-026-groups}{thm:p03-w03-p03-026-groups}{Groups with a single multiplication}
\begin{proof}
(1) A group satisfies associativity, and given \(a,b\) the elements
\(x=a^{-1}b\) and \(y=ba^{-1}\) satisfy \(a\cdot x=b\) and \(y\cdot a=b\).
Conversely let \((G,\cdot)\entails\Gamma\); \(G\) is nonempty.  Fix
\(a\in G\) and choose \(e\) with \(a\cdot e=a\).  For any \(b\) choose \(y\)
with \(y\cdot a=b\); then \(b\cdot e=y\cdot a\cdot e=y\cdot a=b\), so \(e\)
is a right identity for every element.  Choose \(f\) with \(f\cdot a=a\).
For any \(b\) choose \(x\) with \(a\cdot x=b\); then
\(f\cdot b=f\cdot a\cdot x=a\cdot x=b\), so \(f\) is a left identity.
Hence \(f=f\cdot e=e\) is a two-sided identity.  For any \(b\) choose \(x\)
with \(b\cdot x=e\) and \(y\) with \(y\cdot b=e\); then
\(y=y\cdot e=y\cdot b\cdot x=e\cdot x=x\), so \(b\) has a two-sided inverse.
Thus \((G,\cdot)\) is a group.

(2) Both members of \(\Gamma\) are \(\forall\exists\) sentences (the first
with an empty existential block), so (2) follows from (1) and
\cref{thm:p03-w03-p03-026-easy}(3); it is also the implication
(4)\(\Rightarrow\)(1) of \cref{thm:p03-w03-p03-026-chang-los-suszko}.

(3) Write the additive group of integers multiplicatively as
\((\mathbb Z,\cdot)\) with \(m\cdot n=m+n\).  The natural numbers are closed
under this operation, so \((\N,\cdot)\) is a substructure.  It is not a group.
An identity element \(e\) of a group satisfies \(e\cdot0=0\), here \(e+0=0\),
so it would have to be \(0\); and \(1\cdot x=0\) has no solution in \(\N\), so
\(1\) would have no inverse.  So
\(\mathcal G\) is not closed under substructures.  If \(\mathcal G\) were
\(\operatorname{Mod}(U)\) with \(U\) universal, it would be closed under
substructures by \cref{thm:p03-w03-p03-026-easy}(2).
\end{proof}
\end{annexitem}

\subsection{Proofs for \L o\'s 1955}
\label{sec:annex-p03-w03-p03-026-los1955}

In these items \(\mathcal A,\mathcal B,\mathcal C,\mathcal X\) are sets of
formulas, as in \cref{sec:p03-w03-p03-026-los1955}.  \emph{Validity} of a
formula in a model is truth of its universal closure under every assignment,
and \(\mathit{Cn}_i(\mathcal X)\) is read through the handbook's first-order
calculus: \(\alpha\in\mathit{Cn}_i(\mathcal X)\) iff the universal closure of
\(\alpha\) is derivable from the universal closures of the members of
\(\mathcal X\).  This is \L o\'s's operation, because his rule of omitting and
adding general quantifiers makes a formula and its closure interderivable.

\begin{annexitem}{annex:p03-w03-p03-026-los55-cn}{thm:p03-w03-p03-026-los55-cn}{The consequence operation}
\begin{proof}
(1) By soundness and strong completeness
(\cref{thm:fol-soundness,thm:fol-completeness}) for the signature of
\(\mathcal E_i\), the closure of \(\alpha\) is derivable from the closures of
\(\mathcal X\) iff it is true in every structure in which all those closures
are true, that is, iff \(\alpha\) is valid in every model of \(\mathcal X\).
(2.1) is the direction from left to right.  For (2.2), the open formula
\((x_1=x_1)'\) is valid in no model, since domains are nonempty; if
\(\mathcal X\) has a model, \((x_1=x_1)'\notin\mathit{Cn}_i(\mathcal X)\), so
\(\mathcal X\) is consistent.  For G\"odel's theorem, if \(\mathcal X\) is
consistent, some \(\alpha\in\mathcal E_i\) lies outside
\(\mathit{Cn}_i(\mathcal X)\), so some model of \(\mathcal X\) fails to
validate \(\alpha\); in particular \(\mathcal X\) has a model.  Strong
completeness uses the Axiom of Choice exactly when the signs are not
countable, as \cref{thm:fol-completeness} records.

(2) Write \(\operatorname{Mod}(\mathcal X)\) for the models of \(\mathcal X\).
By (1), \(\mathit{Cn}_i(\mathcal X)\) is the set of formulas valid throughout
\(\operatorname{Mod}(\mathcal X)\), so it contains \(\mathcal X\);
\(\operatorname{Mod}(\mathit{Cn}_i(\mathcal X))=\operatorname{Mod}(\mathcal X)\)
by (2.1), whence (1.2).  For (1.3), \(\mathcal Y\subset\mathcal X\) gives
\(\operatorname{Mod}(\mathcal X)\subset\operatorname{Mod}(\mathcal Y)\) and so
\(\mathit{Cn}_i(\mathcal Y)\subset\mathit{Cn}_i(\mathcal X)\).  Conversely,
if \(\alpha\in\mathit{Cn}_i(\mathcal X)\), the closures of \(\mathcal X\)
together with the negation of the closure of \(\alpha\) have no model; by
compactness (\cref{thm:fol-compactness}) some finite part has none, and that
part involves only a finite \(\mathcal Y\subset\mathcal X\), so
\(\alpha\in\mathit{Cn}_i(\mathcal Y)\).  For (1.4), let \(\alpha\) be a
sentence.  For a model \(\mathfrak M\) of \(\mathcal X\), \(\alpha\to\beta\) is
valid in \(\mathfrak M\) iff under every assignment either \(\alpha\) is false
or \(\beta\) is true; since the truth value of \(\alpha\) does not depend on
the assignment, this says: if \(\alpha\) is true in \(\mathfrak M\) then
\(\beta\) is valid in \(\mathfrak M\).  Quantifying over
\(\operatorname{Mod}(\mathcal X)\) gives exactly validity of \(\beta\) in all
models of \(\mathcal X\cup(\alpha)\).  For (1.5) take
\(\alpha=(x_1=x_1)'\): it has no model, so every formula is valid in all its
models.  For (1.6), if \(\mathcal X\) is inconsistent, the \(\alpha\) of
(1.5) lies in \(\mathit{Cn}_i(\mathcal X)\), hence by (1.3) in
\(\mathit{Cn}_i(\mathcal Y)\) for a finite \(\mathcal Y\subset\mathcal X\);
then \(\mathcal Y\) has no model and \(\mathit{Cn}_i(\mathcal Y)=\mathcal E_i\).
For (1.7), if \(\mathcal X\cup(\alpha)\) is inconsistent it has no model, by
(2.2); so \(\alpha\) is false in every model of \(\mathcal X\), and as
\(\alpha\) is a sentence, \(\alpha'\) is valid there.

(3) Let \(g_{t_1},\ldots,g_{t_n}\) be the distinct constants of the sentence
\(\alpha(g_{t_1},\ldots,g_{t_n})\in\mathit{Cn}_3(\mathcal X)\), let
\(\mathfrak N\) be a model for \(\mathcal E_2\) of \(\mathcal X\), and let
\(s\) be an assignment in \(\mathfrak N\).  Expand \(\mathfrak N\) to a model
\(\mathfrak N^+\) for \(\mathcal E_3\) by \(G_{t_k}=s(x_{s_k})\) and arbitrary
values for the other constants; this is well defined because the \(t_k\) are
distinct.  \(\mathcal X\subset\mathcal E_2\) contains no constants, so
\(\mathfrak N^+\) is a model of \(\mathcal X\) (\cref{lem:reduct-invariance})
and \(\alpha(g_{t_1},\ldots,g_{t_n})\) is true in \(\mathfrak N^+\) by (2.1).
The variables \(x_{s_k}\) do not occur in that sentence, so replacing
\(g_{t_k}\) by \(x_{s_k}\) creates no new binding: by induction on terms and
formulas, each term of \(\alpha(x_{s_1},\ldots,x_{s_n})\) takes under \(s\)
(and under any modification of \(s\) at variables other than the
\(x_{s_k}\)) the value of the corresponding term of
\(\alpha(g_{t_1},\ldots,g_{t_n})\) in \(\mathfrak N^+\), and the quantifier
clauses only modify assignments at bound variables, which differ from the
\(x_{s_k}\).  Hence \(\mathfrak N\) satisfies
\(\alpha(x_{s_1},\ldots,x_{s_n})\) under \(s\).  As \(\mathfrak N\) and \(s\)
were arbitrary, \(\alpha(x_{s_1},\ldots,x_{s_n})\) is valid in every model of
\(\mathcal X\), and lies in \(\mathit{Cn}_2(\mathcal X)\) by (1).
\end{proof}
\end{annexitem}

\begin{annexitem}{annex:p03-w03-p03-026-los55-31}{thm:p03-w03-p03-026-los55-31}{Open formulas pass to submodels}
\begin{proof}
Let \(\xi\in\mathcal X\) have free variables \(\bar x\), and let \(\bar a\) be
a tuple from the domain of \(\mathfrak M\).  Since \(\xi\) is valid in
\(\mathfrak M^*\), it is true there at \(\bar a\).  The formula \(\xi\) is
quantifier-free, and \(\equiv\) is a truth function like the other
connectives, so \cref{thm:p03-w03-p03-026-easy}(1) gives that \(\xi\) is
true in \(\mathfrak M\) at \(\bar a\).  Hence \(\xi\) is valid in
\(\mathfrak M\).
\end{proof}
\end{annexitem}

\begin{annexitem}{annex:p03-w03-p03-026-los55-desc}{thm:p03-w03-p03-026-los55-desc}{Descriptions}
\begin{proof}
(4.1) and (4.2): the sentences (i)--(iv) contain no variables at all, so they
have neither free nor bound variables, and their signs are \(=\), the
\(r_k\), the \(f_k\) and constants.

(4.3) Put \(h(a)=G^*_a\).  By (i), \(h\) is injective.  By (iv),
\(F^*_k(h(a_1),\ldots,h(a_{\varkappa_k}))=h(F_k(a_1,\ldots,a_{\varkappa_k}))\),
so \(h[A]\) is closed under the \(F^*_k\) and \(h\) commutes with them.  By
(ii) and (iii), \(R_k(\bar a)\) iff \(R^*_k(h\bar a)\).  So \(h\) is an
isomorphism of \(\mathfrak M\) onto the submodel of \(\mathfrak M^*|1\) with
domain \(h[A]\).  Thus \(\mathfrak M\) is a submodel of \(\mathfrak M^*|1\) in
the generalized sense; since \((\mathfrak M^*|2)|1=\mathfrak M^*|1\), the same
holds for \(\mathfrak M^*|2\) as an extension of the second kind.

(4.4) Let \(\mathfrak M^+\) be the model for \(\mathcal E_3\) obtained from
\(\mathfrak M\) by \(G_a=a\) and, when \(i=1\), arbitrary relations (for
instance empty ones) for the \(q_k\).  Each sentence of
\(\mathcal D(\mathfrak M)\) is true in \(\mathfrak M^+\) by the clause that put
it there.  \(\mathcal X\) contains no constants and, when \(i=1\), no
\(q_k\), so its members are valid in \(\mathfrak M^+\) because they are valid
in \(\mathfrak M\) (\cref{lem:reduct-invariance}).  By (2.2),
\(\mathcal X\cup\mathcal D(\mathfrak M)\) is consistent.

(4.5) By (1.6) some finite subset of
\(\mathcal A\cup\mathcal C\cup\mathcal D(\mathfrak M)\) is inconsistent;
enlarging it keeps it inconsistent, so
\(\mathcal A\cup\mathcal C\cup\mathcal D_0\) is inconsistent for a finite
\(\mathcal D_0\subset\mathcal D(\mathfrak M)\).  If \(\mathcal D_0\) is empty,
\(\mathcal A\cup\mathcal C\) is inconsistent and \(\alpha_1=(x_1=x_1)'\)
serves: it is in \(\mathcal O_1\), in
\(\mathit{Cn}_2(\mathcal A\cup\mathcal C)=\mathcal E_2\), and valid in no model.
Otherwise let \(\alpha=\alpha(g_{a_1},\ldots,g_{a_n})\) be the conjunction of
\(\mathcal D_0\), where \(a_1,\ldots,a_n\) are the distinct elements whose
constants occur.  \(\mathcal A\cup\mathcal C\cup(\alpha)\) has the same models
as \(\mathcal A\cup\mathcal C\cup\mathcal D_0\), so it is inconsistent, and by
(1.7) \(\alpha'\in\mathit{Cn}_3(\mathcal A\cup\mathcal C)\).  As
\(\mathcal A\cup\mathcal C\subset\mathcal E_2\), (1.8) gives
\(\alpha_1=\alpha(x_{s_1},\ldots,x_{s_n})'\in\mathit{Cn}_2(\mathcal A\cup\mathcal C)\)
for distinct variables not occurring in \(\alpha\).  By (4.1) and (4.2),
\(\alpha_1\) has no bound variables, no \(q_k\) and no constants, so
\(\alpha_1\in\mathcal O_1\).  Under the assignment \(x_{s_k}\mapsto a_k\),
every conjunct of \(\alpha(x_{s_1},\ldots,x_{s_n})\) is true in \(\mathfrak M\),
by the definition of \(\mathcal D(\mathfrak M)\), so \(\alpha_1\) is false at
that assignment and is not valid in \(\mathfrak M\).
\end{proof}
\end{annexitem}

\begin{annexitem}{annex:p03-w03-p03-026-los55-thm1}{thm:p03-w03-p03-026-los55-thm1}{The extension theorem}
\begin{proof}
\emph{Necessity.}  Let \(\mathfrak M^*\) be an extension of the second kind of
\(\mathfrak M\) that is a model of \(\mathcal A\cup\mathcal C\).  By (2.1) every
formula of \(\mathit{Cn}_2(\mathcal A\cup\mathcal C)\) is valid in
\(\mathfrak M^*\); a formula of \(\mathcal O_1\) contains only signs of
\(\mathcal E_1\), so it is valid in \(\mathfrak M^*|1\)
(\cref{lem:reduct-invariance}).  \(\mathfrak M\) is a submodel of
\(\mathfrak M^*|1\), so by (3.1) it is a model of
\(\mathcal O_1\cap\mathit{Cn}_2(\mathcal A\cup\mathcal C)\).  If the submodel
relation holds only in the generalized sense, \cref{lem:iso-invariance}
transfers validity along the isomorphism.

\emph{Sufficiency.}  Let \(\mathfrak M\) be a model of
\(\mathcal O_1\cap\mathit{Cn}_2(\mathcal A\cup\mathcal C)\).  If
\(\mathcal A\cup\mathcal C\cup\mathcal D(\mathfrak M)\) were inconsistent in
\(\mathcal E_3\), (4.5) would give a formula of
\(\mathcal O_1\cap\mathit{Cn}_2(\mathcal A\cup\mathcal C)\) not valid in
\(\mathfrak M\).  So the set is consistent, and by G\"odel's theorem it has a
model \(\mathfrak M_0\).  By (4.3), \(h\colon a\mapsto G^{\mathfrak M_0}_a\) is
an isomorphism of \(\mathfrak M\) onto a submodel of \((\mathfrak M_0|2)|1\),
and \(\mathfrak M_0|2\) is a model of \(\mathcal A\cup\mathcal C\), whose
members contain no constants.  To make the extension literal, apply the
copying construction at the start of \cref{sec:annex-p03-w03-p03-026} to
\(h\), with \(\mathfrak M_0|2\) as target; the transport of structure there
carries the relations \(Q_k\) as well.  It yields a model \(\mathfrak M^*\) for
\(\mathcal E_2\) and an isomorphism \(g\colon\mathfrak M^*\to\mathfrak M_0|2\)
extending \(h\), with \(\mathfrak M\) a submodel of \(\mathfrak M^*|1\).  By
\cref{lem:iso-invariance}, \(\mathfrak M^*\) is a model of
\(\mathcal A\cup\mathcal C\).
\end{proof}
\end{annexitem}

\begin{annexitem}{annex:p03-w03-p03-026-los55-remark1}{thm:p03-w03-p03-026-los55-remark1}{Remark~1}
\begin{proof}
Let \(\mathcal X\subset\mathcal E_1\) and \(\alpha\in\mathcal E_1\).  If
\(\alpha\in\mathit{Cn}_1(\mathcal X)\) and \(\mathfrak N\) is a model for
\(\mathcal E_2\) of \(\mathcal X\), then \(\mathfrak N|1\) is a model of
\(\mathcal X\), so \(\alpha\) is valid in \(\mathfrak N|1\) and in
\(\mathfrak N\) (\cref{lem:reduct-invariance}).  If
\(\alpha\in\mathit{Cn}_2(\mathcal X)\) and \(\mathfrak M\) is a model for
\(\mathcal E_1\) of \(\mathcal X\), expand \(\mathfrak M\) by empty relations
for the \(q_k\); the expansion is a model of \(\mathcal X\), so it validates
\(\alpha\), and so does \(\mathfrak M\).  By item~1 of
\cref{thm:p03-w03-p03-026-los55-cn},
\(\mathcal E_1\cap\mathit{Cn}_2(\mathcal X)=\mathit{Cn}_1(\mathcal X)\), and
intersecting with \(\mathcal O_1\) gives the equality.

Now let \(\mathcal C\subset\mathcal E_1\).  If \(\mathfrak M^*\) is an
extension of the first kind of \(\mathfrak M\) and a model of
\(\mathcal A\cup\mathcal C\), its expansion by empty relations for the
\(q_k\) is an extension of the second kind and a model of
\(\mathcal A\cup\mathcal C\); \cref{thm:p03-w03-p03-026-los55-thm1} and the
equality give that \(\mathfrak M\) is a model of
\(\mathcal O_1\cap\mathit{Cn}_1(\mathcal A\cup\mathcal C)\).  Conversely, if
it is, the equality and \cref{thm:p03-w03-p03-026-los55-thm1} give a model
\(\mathfrak M^*\) for \(\mathcal E_2\) of \(\mathcal A\cup\mathcal C\) with
\(\mathfrak M\) a submodel of \(\mathfrak M^*|1\); as
\(\mathcal A\cup\mathcal C\subset\mathcal E_1\), \(\mathfrak M^*|1\) is a model
of \(\mathcal A\cup\mathcal C\) and an extension of the first kind of
\(\mathfrak M\).
\end{proof}
\end{annexitem}

\begin{annexitem}{annex:p03-w03-p03-026-los55-thm2}{thm:p03-w03-p03-026-los55-thm2}{Weak extensions with open secondary conditions}
\begin{proof}
\emph{Necessity.}  Let \(\mathfrak N\) be a model of \(\mathcal A\cup\mathcal C\)
for \(\mathcal E_2\) with \(\mathfrak M\) isomorphic to \(\mathfrak N|1\).  By
(2.1), \(\mathfrak N\) validates
\(\mathcal O_1\cap\mathit{Cn}_2(\mathcal A\cup\mathcal C)\); these formulas
contain only signs of \(\mathcal E_1\), so \(\mathfrak N|1\) validates them,
and so does \(\mathfrak M\) by \cref{lem:iso-invariance}.

\emph{Sufficiency.}  By \cref{thm:p03-w03-p03-026-los55-thm1} there is a model
\(\mathfrak M^*=\langle A^*,R^*_k,F^*_k,Q^*_k\rangle\) of
\(\mathcal A\cup\mathcal C\) with \(\mathfrak M\) a submodel of
\(\mathfrak M^*|1\).  Let \(\mathfrak N\) have domain \(A\), the relations and
functions of \(\mathfrak M\), and \(Q_k=Q^*_k\cap A^{\mu_k}\).  Since
\(\mathfrak M\) is a submodel of \(\mathfrak M^*|1\), \(A\) is closed under the
\(F^*_k\) and \(R^*_k|A=R_k\), \(F^*_k|A=F_k\); with the definition of the
\(Q_k\), \(\mathfrak N\) is a submodel of \(\mathfrak M^*\).  As
\(\mathcal C\subset\mathcal O_2\), (3.1) shows that \(\mathfrak N\) is a model
of \(\mathcal C\).  Also \(\mathfrak N|1=\mathfrak M\), which is a model of
\(\mathcal A\), so \(\mathfrak N\) is a model of \(\mathcal A\)
(\cref{lem:reduct-invariance}).  Thus \(\mathfrak N\) is a weak extension of the
second kind of \(\mathfrak M\), through the identity map, and a model of
\(\mathcal A\cup\mathcal C\).
\end{proof}
\end{annexitem}

\begin{annexitem}{annex:p03-w03-p03-026-los55-thm3}{thm:p03-w03-p03-026-los55-thm3}{\(\mathcal O\)-systems}
\begin{proof}
This is \L o\'s's proof.  Let \(\alpha\in\mathcal X\) and suppose
\(\alpha\notin\mathit{Cn}_1(\mathcal O_1\cap\mathcal X)\).  By item~1 of
\cref{thm:p03-w03-p03-026-los55-cn} some model \(\mathfrak M\) of
\(\mathcal O_1\cap\mathcal X\) does not validate \(\alpha\).  Put
\(\mathcal A=\mathcal O_1\cap\mathcal X\) and \(\mathcal C=\mathcal X\).  Then
\(\mathcal A\cup\mathcal C=\mathcal X\subset\mathcal E_1\), and by
\cref{thm:p03-w03-p03-026-los55-remark1} and \(\mathit{Cn}_1(\mathcal X)=\mathcal X\),
\[
 \mathcal O_1\cap\mathit{Cn}_2(\mathcal A\cup\mathcal C)
 =\mathcal O_1\cap\mathit{Cn}_1(\mathcal X)=\mathcal O_1\cap\mathcal X=\mathcal A ,
\]
of which \(\mathfrak M\) is a model.  By
\cref{thm:p03-w03-p03-026-los55-remark1}, \(\mathfrak M\) has an extension of
the first kind \(\mathfrak M^*\) that is a model of \(\mathcal X\).  By
hypothesis the submodel \(\mathfrak M\) of \(\mathfrak M^*\) is a model of
\(\mathcal X\), so it validates \(\alpha\), a contradiction.  Hence
\(\mathcal X\subset\mathit{Cn}_1(\mathcal O_1\cap\mathcal X)\).  Conversely,
\(\mathit{Cn}_1(\mathcal O_1\cap\mathcal X)\subset\mathit{Cn}_1(\mathcal X)=\mathcal X\)
by monotonicity (\cref{annex:p03-w03-p03-026-los55-cn}, (1.3)).  So
\(\mathcal X=\mathit{Cn}_1(\mathcal O_1\cap\mathcal X)\).
\end{proof}
\end{annexitem}

\begin{annexitem}{annex:p03-w03-p03-026-los55-thm45}{thm:p03-w03-p03-026-los55-thm45}{Theorems~4 and~5}
\begin{proof}
\emph{Theorem~5.}  Let \(\mathbf B=\mathbf M(\mathcal B')\) with
\(\mathcal B'\subset\mathcal E_1\).  As \(\mathbf B\subset\mathbf A\),
\(\mathbf B=\mathbf M(\mathcal A)\cap\mathbf M(\mathcal B')=\mathbf M(\mathcal A\cup\mathcal B')\).
Put \(\mathcal B=\mathcal O_1\cap\mathit{Cn}_1(\mathcal A\cup\mathcal B')\subset\mathcal O_1\).
Every member of \(\mathbf B\) is a model of \(\mathcal A\cup\mathcal B'\), hence
of \(\mathit{Cn}_1(\mathcal A\cup\mathcal B')\supset\mathcal B\) by (2.1); so
\(\mathbf B\subset\mathbf M(\mathcal A\cup\mathcal B)\).  Conversely let
\(\mathfrak M\) be a model of \(\mathcal A\cup\mathcal B\).  Apply
\cref{thm:p03-w03-p03-026-los55-remark1} with primary conditions
\(\mathcal A\) and secondary conditions \(\mathcal B'\subset\mathcal E_1\): its
condition is that \(\mathfrak M\) be a model of
\(\mathcal O_1\cap\mathit{Cn}_1(\mathcal A\cup\mathcal B')=\mathcal B\), which
holds.  So \(\mathfrak M\) has an extension of the first kind \(\mathfrak M^*\)
that is a model of \(\mathcal A\cup\mathcal B'\), that is, \(\mathfrak M^*\in\mathbf B\).
Since \(\mathfrak M\) is a model of \(\mathcal A\), it is an
\(\mathbf A\)-submodel of \(\mathfrak M^*\), and the hypothesis puts it in
\(\mathbf B\).  Hence \(\mathbf B=\mathbf M(\mathcal A\cup\mathcal B)\).

\emph{Theorem~4.}  Let \(\mathbf B\) be the given class and apply Theorem~5
with \(\mathcal A=\varnothing\), so that \(\mathbf A=\mathbf M(\varnothing)\) is
the class of all models for \(\mathcal E_1\) and the \(\mathbf A\)-submodels
of a model are all its submodels.  The hypothesis of Theorem~5 is then the
hypothesis of Theorem~4, and it yields \(\mathcal B\subset\mathcal O_1\) with
\(\mathbf B=\mathbf M(\mathcal B)\).
\end{proof}
\end{annexitem}

\begin{annexitem}{annex:p03-w03-p03-026-los55-81}{thm:p03-w03-p03-026-los55-81}{(8.1)}
\begin{proof}
The class \(\mathbf M(\mathcal A\cup(\alpha))\) is a definable subclass of
\(\mathbf A\), and the hypothesis is that of Theorem~5 for it.  So
\(\mathbf M(\mathcal A\cup(\alpha))=\mathbf M(\mathcal A\cup\mathcal B)\) for
some \(\mathcal B\subset\mathcal O_1\).  Then
\(\alpha\in\mathit{Cn}_1(\mathcal A\cup\mathcal B)\) by item~1 of
\cref{thm:p03-w03-p03-026-los55-cn}, and by (1.3)
\(\alpha\in\mathit{Cn}_1(\mathcal A\cup\mathcal B_0)\) for a finite
\(\mathcal B_0=\{\gamma_1,\ldots,\gamma_k\}\subset\mathcal B\).  Let
\(\gamma=\gamma_1\land\cdots\land\gamma_k\) (with \(\gamma\) the formula
\(x_1=x_1\) if \(k=0\)), an open formula, and let \(\beta\) be
\(\prod_{x_{s_1}}\cdots\prod_{x_{s_n}}\gamma\), where
\(x_{s_1},\ldots,x_{s_n}\) are the free variables of \(\gamma\).  A model
validates \(\beta\) iff it validates every \(\gamma_j\), because validity is
truth under all assignments and a conjunction is true iff its conjuncts are.
Hence \(\mathcal A\cup(\beta)\) and \(\mathcal A\cup\mathcal B_0\) have the same
models, \(\alpha\in\mathit{Cn}_1(\mathcal A\cup(\beta))\), and
\(\beta\to\alpha\in\mathit{Cn}_1(\mathcal A)\) by (1.4).  In the other
direction, every \(\gamma_j\in\mathcal B\) is valid in all models of
\(\mathcal A\cup\mathcal B\), which are the models of
\(\mathcal A\cup(\alpha)\); so \(\beta\in\mathit{Cn}_1(\mathcal A\cup(\alpha))\)
and \(\alpha\to\beta\in\mathit{Cn}_1(\mathcal A)\) by (1.4).  Both being
sentences, \(\alpha\equiv\beta\) is valid in every model of \(\mathcal A\), so
it belongs to \(\mathit{Cn}_1(\mathcal A)\).
\end{proof}
\end{annexitem}

\begin{annexitem}{annex:p03-w03-p03-026-los55-thm6}{thm:p03-w03-p03-026-los55-thm6}{Persistent sentences}
\begin{proof}
We check the hypothesis of (8.1) for the sentence \(\alpha'\).  Let
\(\mathfrak N\in\mathbf M(\mathcal A\cup(\alpha'))\) and let \(\mathfrak M\) be
an \(\mathbf A\)-submodel of \(\mathfrak N\).  If \(\alpha\) were valid in
\(\mathfrak M\), then, since \(\mathfrak N\in\mathbf A\) is an extension of
\(\mathfrak M\), persistence would make \(\alpha\) valid in \(\mathfrak N\),
where \(\alpha'\) is valid.  So \(\alpha\) is not valid in \(\mathfrak M\), and
as it is a sentence, \(\alpha'\) is; thus
\(\mathfrak M\in\mathbf M(\mathcal A\cup(\alpha'))\).  By (8.1),
\(\alpha'\equiv\prod_{x_{s_1}}\cdots\prod_{x_{s_n}}\gamma\) belongs to
\(\mathit{Cn}_1(\mathcal A)\) for some \(\gamma\in\mathcal O_1\).  In every
model, \(\bigl(\prod_{x_{s_1}}\cdots\prod_{x_{s_n}}\gamma\bigr)'\) has the
truth value of \(\sum_{x_{s_1}}\cdots\sum_{x_{s_n}}\gamma'\)
(\cref{thm:p03-w03-p03-026-prenex}), and \(\gamma'\in\mathcal O_1\).  Hence
\(\alpha\equiv\sum_{x_{s_1}}\cdots\sum_{x_{s_n}}\gamma'\) is valid in all
models of \(\mathcal A\) and belongs to \(\mathit{Cn}_1(\mathcal A)\).
\end{proof}
\end{annexitem}

\begin{annexitem}{annex:p03-w03-p03-026-los55-thm789}{thm:p03-w03-p03-026-los55-thm789}{Theorems~7, 8 and~9}
\begin{proof}
\emph{Theorem~8.}  A model \(\mathfrak M\) for \(\mathcal E_1\) lies in
\(\operatorname{sm}(\mathbf A|1)\) iff it is isomorphic to a submodel of
\(\mathfrak N|1\) for some model \(\mathfrak N\) of \(\mathcal A\).  By the
copying construction of \cref{sec:annex-p03-w03-p03-026}, applied as in
\cref{annex:p03-w03-p03-026-los55-thm1} with the \(Q_k\) transported, and by
\cref{lem:iso-invariance}, this holds iff \(\mathfrak M\) has an extension of
the second kind that is a model of \(\mathcal A\): an extension of
\(\mathfrak M\) with primary conditions \(\varnothing\) and secondary
conditions \(\mathcal A\).  By \cref{thm:p03-w03-p03-026-los55-thm1} this holds
iff \(\mathfrak M\) is a model of
\(\mathcal O_1\cap\mathit{Cn}_2(\mathcal A)\).

\emph{Theorem~7.}  Let \(\mathcal A\subset\mathcal E_1\) and let
\(\mathbf A_2\) be the class of models for \(\mathcal E_2\) of \(\mathcal A\).
Every model for \(\mathcal E_1\) of \(\mathcal A\) is the reduct of its
expansion by empty relations for the \(q_k\), and every reduct of a member of
\(\mathbf A_2\) is a model of \(\mathcal A\); so \(\mathbf A_2|1=\mathbf A\).
Theorem~8 gives
\(\operatorname{sm}(\mathbf A)=\mathbf M(\mathcal O_1\cap\mathit{Cn}_2(\mathcal A))\),
and \cref{thm:p03-w03-p03-026-los55-remark1} replaces \(\mathit{Cn}_2\) by
\(\mathit{Cn}_1\).

\emph{Theorem~9.}  By Theorem~8,
\(\mathbf B\cap\operatorname{sm}(\mathbf A|1)=\mathbf M(\mathcal B)\cap\mathbf M(\mathcal O_1\cap\mathit{Cn}_2(\mathcal A))
=\mathbf M[\mathcal B\cup(\mathcal O_1\cap\mathit{Cn}_2(\mathcal A))]\).  If
\(\mathfrak N\) is a model of \(\mathcal A\), then \(\mathfrak N|1\in\mathbf B\)
validates \(\mathcal B\), and so does \(\mathfrak N\); hence
\(\mathcal B\subset\mathit{Cn}_2(\mathcal A)\), and by (1.2)
\(\mathit{Cn}_2(\mathcal A)\subset\mathit{Cn}_2(\mathcal A\cup\mathcal B)\subset\mathit{Cn}_2(\mathit{Cn}_2(\mathcal A))=\mathit{Cn}_2(\mathcal A)\).
So \(\mathit{Cn}_2(\mathcal A\cup\mathcal B)=\mathit{Cn}_2(\mathcal A)\), which
gives the stated form.
\end{proof}
\end{annexitem}

\begin{annexitem}{annex:p03-w03-p03-026-los55-thm1011}{thm:p03-w03-p03-026-los55-thm1011}{Theorems~10 and~11}
\begin{proof}
\emph{(11.1).}  If \(\mathfrak M\in\mathbf M(\mathcal A,\alpha)\), choose
\(Q_1,\ldots,Q_s\) making \(\beta\) valid and interpret the \(q_j\), \(j>s\),
arbitrarily; the result is a model for \(\mathcal E_2\) of
\(\mathcal A\cup(\beta)\), since \(\beta\) contains no \(q_j\) with \(j>s\)
and \(\mathcal A\) no \(q_j\) at all, and its reduct is \(\mathfrak M\).
Conversely, the reduct of a model of \(\mathcal A\cup(\beta)\) is a model of
\(\mathcal A\) in which its own \(Q_1,\ldots,Q_s\) witness \(\alpha\).

\emph{Theorem~10.}  Here \(\beta=\prod_{x_1}\cdots\prod_{x_n}\gamma\) is a
sentence, so every free variable of \(\gamma\) is among
\(x_1,\ldots,x_n\), and a model for \(\mathcal E_2\) validates \(\gamma\) iff
it validates \(\beta\) (quantifiers over variables not free in \(\gamma\)
are vacuous, domains being nonempty).  Hence
\(\mathit{Cn}_2(\mathcal A\cup(\gamma))=\mathit{Cn}_2(\mathcal A\cup(\beta))\).
Let \(\mathfrak M\) be a model of \(\mathcal A\).  By (11.1),
\(\mathfrak M\in\mathbf M(\mathcal A,\alpha)\) iff \(\mathfrak M=\mathfrak N|1\)
for a model \(\mathfrak N\) of \(\mathcal A\cup(\gamma)\).  This holds iff
\(\mathfrak M\) has a weak extension of the second kind that is a model of
\(\mathcal A\cup(\gamma)\): given one, \(\mathfrak N\) with an isomorphism
\(\varphi\colon\mathfrak M\to\mathfrak N|1\), pull the relations
\(Q_k^{\mathfrak N}\) back along \(\varphi\) to get an expansion of
\(\mathfrak M\) isomorphic to \(\mathfrak N\), which is a model of
\(\mathcal A\cup(\gamma)\) by \cref{lem:iso-invariance}.  As
\((\gamma)\subset\mathcal O_2\), \cref{thm:p03-w03-p03-026-los55-thm2} says
that such a weak extension exists iff \(\mathfrak M\) is a model of
\(\mathcal O_1\cap\mathit{Cn}_2(\mathcal A\cup(\gamma))\).  Therefore
\(\mathbf M(\mathcal A,\alpha)=\mathbf M\bigl(\mathcal A\cup(\mathcal O_1\cap\mathit{Cn}_2(\mathcal A\cup(\beta)))\bigr)\).

\emph{Theorem~11.}  By (11.1),
\(\mathbf M(\mathcal A,\alpha)=\mathbf A_\beta|1\), where \(\mathbf A_\beta\) is
the class of models for \(\mathcal E_2\) of
\(\mathcal A\cup(\beta)\subset\mathcal E_2\).  Theorem~8 gives
\(\operatorname{sm}(\mathbf M(\mathcal A,\alpha))=\mathbf M(\mathcal O_1\cap\mathit{Cn}_2(\mathcal A\cup(\beta)))\).
Since \(\mathbf A_\beta|1\subset\mathbf M(\mathcal A)\), Theorem~9 with
\(\mathbf B=\mathbf M(\mathcal A)\) gives
\(\mathbf M(\mathcal A)\cap\operatorname{sm}(\mathbf M(\mathcal A,\alpha))
=\mathbf M\bigl(\mathcal A\cup(\mathcal O_1\cap\mathit{Cn}_2(\mathcal A\cup(\beta)\cup\mathcal A))\bigr)\),
and \(\mathcal A\cup(\beta)\cup\mathcal A=\mathcal A\cup(\beta)\).
\end{proof}
\end{annexitem}

\begin{annexitem}{annex:p03-w03-p03-026-los55-arithmetic}{thm:p03-w03-p03-026-los55-arithmetic}{The arithmetic example}
\begin{proof}
Write \(\mathbb N\) for the standard model.  Define numerals by
\(\bar 0=f_3\) and \(\overline{n+1}=f_1(\bar n,f_4)\), and for a model
\(\mathfrak M\) of \(\mathcal A\) let \(\nu(n)\) be the value of \(\bar n\) in
\(\mathfrak M\).

(a) \(\nu\) embeds \(\mathbb N\) in \(\mathfrak M\).  The set \(\mathcal A\)
contains \((\bar n=\bar m)'\) for \(n\neq m\),
\(f_1(\bar n,\bar m)=\overline{n+m}\), \(f_2(\bar n,\bar m)=\overline{nm}\),
\(f_3=\bar 0\) and \(f_4=\bar 1\), all valid in \(\mathbb N\).  So \(\nu\) is
injective and commutes with the operations and constants.

(b) \(\alpha\) is valid in a model \(\mathfrak M\) of \(\mathcal A\) iff \(\nu\)
is not onto.  If \(\nu\) is not onto, let \(Q_1=\nu[\mathbb N]\): it contains
the value of \(f_3\), which is \(\nu(0)\); if it contains \(\nu(n)\), it
contains \(f_1(\nu(n),f_4)=\nu(n+1)\) by (a); and it misses some element.  So
the three conjuncts hold.  If \(\nu\) is onto and \(Q_1\) contains
\(\nu(0)\) and is closed under \(x\mapsto f_1(x,f_4)\), then
\(\nu(n)\in Q_1\) for every \(n\) by induction on \(n\), so \(Q_1\) is
everything and the third conjunct fails.  The same argument shows that the
induction axiom (every set containing \(0\) and closed under successor is
everything) holds in \(\mathfrak M\) iff \(\nu\) is onto, that is, iff
\(\mathfrak M\) is isomorphic to \(\mathbb N\).  So
\(\mathbf M(\mathcal A,\alpha)\) is the class of models of \(\mathcal A\) in
which induction fails.

(c) \(\mathbb N\in\mathbf M(\mathcal A)\setminus\mathbf M(\mathcal A,\alpha)\),
since there \(\nu\) is the identity.

(d) \(\mathbf M(\mathcal A,\alpha)\neq\varnothing\).  Add one constant \(g\)
and let \(\Sigma=\mathcal A\cup\{(g=\bar n)':n\in\mathbb N\}\).  A finite part
of \(\Sigma\) mentions finitely many \(n\) and is valid in \(\mathbb N\) with
\(g\) interpreted as a number larger than all of them.  By compactness
(\cref{thm:fol-compactness}) \(\Sigma\) has a model; in its reduct, which is a
model of \(\mathcal A\), the value of \(g\) is not of the form \(\nu(n)\), so
by (b) the reduct lies in \(\mathbf M(\mathcal A,\alpha)\).

(e) Suppose \(\mathbf M(\mathcal A,\alpha)=\mathbf M(\mathcal B)\) for some
\(\mathcal B\subset\mathcal E_1\).  Choose \(\mathfrak N\) in the class by (d).
For \(\delta\in\mathcal B\) let \(\bar\delta\) be its closure.  Either
\(\bar\delta\) or \(\bar\delta'\) is valid in \(\mathbb N\), so one of them is
in \(\mathcal A\).  \(\mathfrak N\) is a model of \(\mathcal A\) and validates
\(\bar\delta\), so \(\bar\delta'\notin\mathcal A\), hence
\(\bar\delta\in\mathcal A\) and \(\delta\) is valid in \(\mathbb N\).  Thus
\(\mathbb N\in\mathbf M(\mathcal B)=\mathbf M(\mathcal A,\alpha)\),
contradicting (c).  So \(\mathbf M(\mathcal A,\alpha)\) is not elementarily
definable.
\end{proof}
\end{annexitem}

\begin{annexitem}{annex:p03-w03-p03-026-los55-semigroup}{thm:p03-w03-p03-026-los55-semigroup}{Semigroups extending to groups}
\begin{proof}
The models of \(\mathcal A_1\cup\mathcal C_1\) are exactly the groups: the
set \(\mathcal A_1\cup\mathcal C_1\) has the same models as the two sentences of
\cref{thm:p03-w03-p03-026-groups}(1), because
\(\prod_{a}\prod_{b}\sum_{x}\sum_{y}(a\cdot x=b\land y\cdot a=b)\) holds iff
both members of \(\mathcal C_1\) hold.  As \(\mathcal C_1\subset\mathcal E_1\),
\cref{thm:p03-w03-p03-026-los55-remark1} with primary conditions
\(\mathcal A_1\) and secondary conditions \(\mathcal C_1\) gives the
equivalence.  The only sign of \(\mathcal E_1\) is \(\cdot\), so the formulas
of \(\mathcal O_1\) contain no other operation or relation sign.
\end{proof}
\end{annexitem}

\begin{annexitem}{annex:p03-w03-p03-026-los55-ordered}{thm:p03-w03-p03-026-los55-ordered}{Orderable groups}
\begin{proof}
Let \(G\) be a group and \(Q\) a binary relation on it.  \(\gamma^+\) is
valid under \(q_1\mapsto Q\) iff: \(Q(a,b)\) excludes \(Q(b,a)\), which with
\(a=b\) makes \(Q\) irreflexive; \(Q\) is transitive; any two distinct
elements are comparable; and \(Q(a,b)\) implies \(Q(cad,cbd)\).  These say
exactly that \(Q\) is an ordering relation \(<\) of the group in the sense
of \S\,15.  So \(\alpha^+\) is valid in \(G\) iff \(G\) is orderable, and the
class of orderable groups is \(\mathbf M(\mathcal A_1\cup\mathcal C_1,\alpha^+)\),
the models of \(\mathcal A_1\cup\mathcal C_1\) being the groups
(\cref{thm:p03-w03-p03-026-los55-semigroup}).  The formula \(\gamma^+\) is
open and lies in \(\mathcal O_2\), with \(q_1\) the only new sign, so
\(\beta^+\) is a general sentence and Theorem~10
(\cref{thm:p03-w03-p03-026-los55-thm1011}) gives the stated axioms.
\end{proof}
\end{annexitem}

\begin{annexitem}{annex:p03-w03-p03-026-los55-ternary}{thm:p03-w03-p03-026-los55-ternary}{A reducible extension of a ternary semigroup}
\begin{proof}
Let \(T\) be a ternary semigroup; write its operation \([a,b,c]\), so that
complete associativity reads
\([[a,b,c],d,e]=[a,[b,c,d],e]=[a,b,[c,d,e]]\).  On \(T\times T\) put
\((a,b)\sim(a',b')\) iff \([a,b,x]=[a',b',x]\) and \([x,a,b]=[x,a',b']\) for
every \(x\in T\); this is an equivalence relation.  Let \(P\) be the set of
classes \(\langle a,b\rangle\), taken disjoint from \(T\) (replace it by a
copy if necessary), and \(S=T\cup P\).  Define, for \(a,b,c,d\in T\),
\begin{gather*}
 a\cdot b=\langle a,b\rangle,\qquad
 a\cdot\langle b,c\rangle=[a,b,c],\\
 \langle a,b\rangle\cdot c=[a,b,c],\qquad
 \langle a,b\rangle\cdot\langle c,d\rangle=\langle[a,b,c],d\rangle .
\end{gather*}
\emph{Well defined.}  If \((b,c)\sim(b',c')\), then
\([a,b,c]=[a,b',c']\) by the second clause of \(\sim\) with \(x=a\); the
third product is handled by the first clause.  In the fourth, replacing
\((a,b)\) by \((a',b')\sim(a,b)\) leaves \([a,b,c]\) unchanged.  Replacing
\((c,d)\) by \((c',d')\sim(c,d)\): for every \(x\),
\begin{align*}
 [[a,b,c],d,x]&=[a,b,[c,d,x]]=[a,b,[c',d',x]]=[[a,b,c'],d',x],\\
 [x,[a,b,c],d]&=[[x,a,b],c,d]=[[x,a,b],c',d']=[x,[a,b,c'],d'],
\end{align*}
using complete associativity and the two clauses of \(\sim\).  So
\(([a,b,c],d)\sim([a,b,c'],d')\).

\emph{Associative.}  There are eight cases, by whether each of three factors
\(u,v,w\) lies in \(T\) or in \(P\).  With \(a,\ldots,f\in T\):
\begin{itemize}
 \item \(T,T,T\): \((a\cdot b)\cdot c=[a,b,c]=a\cdot(b\cdot c)\).
 \item \(T,T,P\) and \(T,P,T\): both sides reduce to the pair classes of
       \(([a,b,c],d)\) and \((a,[b,c,d])\), namely
       \((a\cdot b)\cdot\langle c,d\rangle=\langle[a,b,c],d\rangle\),
       \(a\cdot(b\cdot\langle c,d\rangle)=\langle a,[b,c,d]\rangle\),
       \((a\cdot\langle b,c\rangle)\cdot d=\langle[a,b,c],d\rangle\),
       \(a\cdot(\langle b,c\rangle\cdot d)=\langle a,[b,c,d]\rangle\).  These
       classes coincide, because for every \(x\)
       \begin{align*}
        [[a,b,c],d,x]&=[a,[b,c,d],x],\\
        [x,[a,b,c],d]&=[[x,a,b],c,d]=[x,a,[b,c,d]].
       \end{align*}
 \item \(P,T,T\): \((\langle a,b\rangle\cdot c)\cdot d=\langle[a,b,c],d\rangle\), and
       \(\langle a,b\rangle\cdot(c\cdot d)=\langle a,b\rangle\cdot\langle c,d\rangle=\langle[a,b,c],d\rangle\).
 \item \(T,P,P\): \((a\cdot\langle b,c\rangle)\cdot\langle d,e\rangle=[[a,b,c],d,e]\)
       and \(a\cdot(\langle b,c\rangle\cdot\langle d,e\rangle)=[a,[b,c,d],e]\).
 \item \(P,T,P\): \((\langle a,b\rangle\cdot c)\cdot\langle d,e\rangle=[[a,b,c],d,e]\)
       and \(\langle a,b\rangle\cdot(c\cdot\langle d,e\rangle)=[a,b,[c,d,e]]\).
 \item \(P,P,T\): \((\langle a,b\rangle\cdot\langle c,d\rangle)\cdot e=[[a,b,c],d,e]\)
       and \(\langle a,b\rangle\cdot(\langle c,d\rangle\cdot e)=[a,b,[c,d,e]]\).
 \item \(P,P,P\):
       \((\langle a,b\rangle\cdot\langle c,d\rangle)\cdot\langle e,f\rangle=\langle[[a,b,c],d,e],f\rangle\)
       and
       \(\langle a,b\rangle\cdot(\langle c,d\rangle\cdot\langle e,f\rangle)=\langle[a,b,[c,d,e]],f\rangle\).
\end{itemize}
In the last four cases the two sides agree by complete associativity.

\emph{Conclusions.}  Put \((u,v,w)=(u\cdot v)\cdot w\) on \(S\).  Then
\(\langle S,\cdot,(\,,\,,\,)\rangle\) is a model of
\(\mathcal A_3\cup\mathcal C_3\): \(\mathcal C_3\) holds by associativity and
the definition, and each member of \(\mathcal A_3\) holds because its sides are
products of the same five factors in the same order.  For \(a,b,c\in T\),
\((a\cdot b)\cdot c=\langle a,b\rangle\cdot c=[a,b,c]\in T\), so \(T\) is closed
under the ternary operation of \(S\), which restricts to that of \(T\): \(T\)
is a submodel of the reduct of \(S\) to \(\mathcal E_1\).  Hence
\(\mathbf T\subset\operatorname{sm}(\mathbf M(\mathcal A_3\cup\mathcal C_3)|1)\).
Conversely, the reduct of a model of \(\mathcal A_3\cup\mathcal C_3\) is a
model of \(\mathcal A_3\subset\mathcal O_1\), and by (3.1) and
\cref{lem:iso-invariance} so is each submodel in the generalized sense; so the
classes are equal.  Finally,
\(\mathcal O_1\cap\mathit{Cn}_1(\mathcal A_3)\subset\mathcal O_1\cap\mathit{Cn}_2(\mathcal A_3\cup\mathcal C_3)\)
by \cref{thm:p03-w03-p03-026-los55-remark1} and monotonicity.  If
\(\xi\in\mathcal O_1\cap\mathit{Cn}_2(\mathcal A_3\cup\mathcal C_3)\) and \(T\)
is a ternary semigroup, the model \(S\) above validates \(\xi\) by (2.1), so
its reduct does, and so does its submodel \(T\) by (3.1); by item~1 of
\cref{thm:p03-w03-p03-026-los55-cn}, \(\xi\in\mathit{Cn}_1(\mathcal A_3)\).
\end{proof}
\end{annexitem}

\begin{annexitem}{annex:p03-w03-p03-026-los55-omega}{thm:p03-w03-p03-026-los55-omega}{The proved cases of the \S\,17 theorem}
\begin{proof}
(1) Let \(\mathfrak N\) be a model for \(\mathcal E_2\) of
\(\mathcal A\cup\mathcal C\) and \(\varphi\colon\mathfrak M\to\mathfrak N|1\) an
isomorphism.  Expand \(\mathfrak N\) to \(\mathfrak N^+\) by
\(G_a=\varphi(a)\), and let \(\mathcal X\) be the set of formulas of
\(\mathcal E_3\) valid in \(\mathfrak N^+\).  By item~1 of
\cref{thm:p03-w03-p03-026-los55-cn}, \(\mathcal X\) is a system, and it is
consistent by (2.2).  It is complete, since a sentence or its negation is true
in \(\mathfrak N^+\).  It is \(\omega\)-complete: if \(\alpha(g_a)\in\mathcal X\)
for every \(a\in A\), then, as in the proof of (1.8), \(\alpha(x_s)\) holds in
\(\mathfrak N^+\) when \(x_s\) takes the value \(\varphi(a)\), and \(\varphi\) is
onto, so \(\prod_{x_s}\alpha(x_s)\) is true.  It contains
\(\mathcal A\cup\mathcal C\) (\cref{lem:reduct-invariance}) and
\(\mathcal D(\mathfrak M)\), because \(\varphi\) is an isomorphism.

(2) Let \(\mathcal X\) be as in item~2, and \(\mathfrak N\) a model of
\(\mathcal X\), which exists by G\"odel's theorem.  It suffices to find a
model \(\mathfrak N_0\) of \(\mathcal X\) every element of which is some
\(G_a^{\mathfrak N_0}\).  Indeed, the map \(a\mapsto G_a^{\mathfrak N_0}\) is then
onto, injective by (i), and an isomorphism of \(\mathfrak M\) onto
\(\mathfrak N_0|1\) by (ii)--(iv); so \(\mathfrak N_0|2\) is a weak extension of
the second kind of \(\mathfrak M\), and it is a model of
\(\mathcal A\cup\mathcal C\subset\mathcal X\).

\emph{\(\mathcal X\) complete.}  Let \(N_0=\{G_a^{\mathfrak N}:a\in A\}\).  It
contains the values of all constants, and by (iv) it is closed under the
operations, so it is the domain of a substructure \(\mathfrak N_0\) of
\(\mathfrak N\).  We verify the Tarski--Vaught condition
(\cref{thm:tarski-vaught}).  Every element of \(N_0\) is named by a constant,
so it suffices to treat a sentence \(\sum_{x}\varphi(x)\) of \(\mathcal E_3\)
true in \(\mathfrak N\).  Its negation is false in \(\mathfrak N\), so not in
\(\mathcal X\); by completeness \(\sum_x\varphi(x)\in\mathcal X\).  If
\(\varphi(g_a)\notin\mathcal X\) for every \(a\), completeness gives
\(\varphi(g_a)'\in\mathcal X\) for every \(a\), and \(\omega\)-completeness
gives \(\prod_x\varphi(x)'\in\mathcal X\), which contradicts consistency
together with \(\sum_x\varphi(x)\in\mathcal X\).  So
\(\varphi(g_a)\in\mathcal X\) for some \(a\), and \(\mathfrak N\) satisfies
\(\varphi\) at \(G_a^{\mathfrak N}\in N_0\).  Hence
\(\mathfrak N_0\preccurlyeq\mathfrak N\), and \(\mathfrak N_0\) is a model of
\(\mathcal X\).

\emph{\(A\) and the signs countable.}  Then \(\mathcal E_3\) has a countable
signature.  Let \(\Sigma\) be the set of universal closures of members of
\(\mathcal X\), a satisfiable theory, and
\(p(x)=\{(x=g_a)':a\in A\}\).  If \(p\) is not consistent with \(\Sigma\),
every model of \(\Sigma\) omits it.  Otherwise \(p\) is nonprincipal over
\(\Sigma\) (\cref{def:principal-type}): suppose \(\theta(x)\) is a formula with
\(\Sigma\cup\{\sum_x\theta(x)\}\) satisfiable and
\(\Sigma\entails\prod_x(\theta(x)\to(x=g_a)')\) for every \(a\).  Then
\(\theta(g_a)'\) is valid in every model of \(\mathcal X\), so
\(\theta(g_a)'\in\mathit{Cn}_3(\mathcal X)=\mathcal X\) for every \(a\), and
\(\omega\)-completeness gives \(\prod_x\theta(x)'\in\mathcal X\), contradicting
satisfiability of \(\Sigma\cup\{\sum_x\theta(x)\}\).  By
\cref{thm:p03-vaught-partial-omission}, \(\Sigma\) has a model
\(\mathfrak N_0\) omitting \(p\), that is, a model of \(\mathcal X\) every
element of which is some \(G_a^{\mathfrak N_0}\).
\end{proof}
\end{annexitem}
